\clearpage
\chapter{Literature Review}
\section{Review scope and selection}
The review covers twelve primary research works. The search used Google Scholar, PubMed, Europe PMC, IEEE Xplore and the Nature and Frontiers portals with combinations of the keywords \emph{glioma segmentation}, \emph{brain tumour segmentation}, \emph{BraTS}, \emph{multi-modal MRI}, \emph{transformer}, \emph{2.5D}, \emph{cross-slice} and \emph{inter-slice context}. Works were selected if they were peer-reviewed, addressed glioma analysis on multi-modal MRI or introduced a segmentation architecture used as a reference in this field, and reported a quantitative evaluation. The selection spans 2015 to 2026 and has two parts. The first part consists of seven recent works (2024--2026): five on glioma analysis, which define the current state of the application area, our previous-semester reference paper, and the cross-slice transformer that is the base paper of this project. The second part consists of five foundational architectures (2015--2022) against which glioma segmentation methods are routinely compared. Bibliographic details were verified against the publishers' records; quantitative findings are taken from the papers' abstracts or full texts, and details that we could not verify are marked ``Not reported''.

\section{Descriptive review of literature works}

\subsection{MM-MSCA-AF: multi-scale contextual aggregation and attention fusion}
Aslam \emph{et al.} (2025), in ``Enhanced brain tumor segmentation in medical imaging using multi-modal multi-scale contextual aggregation and attention fusion'' \cite{aslam2025}, investigate brain-tumour segmentation from multi-modal MRI, where tumour heterogeneity and variability across sequences make segmentation difficult. The work proposes the MM-MSCA-AF framework, which takes T1, T2, FLAIR and T1ce images, uses multi-scale contextual aggregation to obtain both global and fine-grained spatial features, and applies gated attention fusion to keep informative features and suppress noise. The study is set on the BraTS 2020 dataset. The authors evaluate the model with the Dice coefficient and report a Dice of 0.8158 for the necrotic tumour region and 0.8589 overall, outperforming U-Net, nnU-Net and Attention U-Net in their comparison. Limitations stated beyond the single-dataset evaluation are not reported in the abstract, and boundary metrics such as HD95 are not reported there. The model aggregates context within a slice or patch and does not explicitly model the relationship between neighbouring slices; this gap, together with its use of the same dataset as our project, makes it a direct point of comparison.

\subsection{DeepGlioSeg: CNN--transformer parallel combination}
Li \emph{et al.} (2025), in ``DeepGlioSeg: advanced glioma MRI data segmentation with integrated local-global representation architecture'' \cite{li2025deepglioseg}, address the heterogeneity of gliomas and the imbalance between tumour regions, which hinder existing segmentation methods. They propose a U-shaped network whose core is a CNN--Transformer Parallel Combination (CTPC) module: parallel convolutional and transformer branches fuse local and global features. A region-based probability, computed from the ratio of tumour to background pixels, gives more weight to under-represented regions, and test-time augmentation with volume-constrained post-processing refines the output. The study uses three public glioma MRI datasets and one private clinical dataset. The authors report, from quantitative and qualitative comparisons, that DeepGlioSeg outperforms other state-of-the-art methods; the per-region values were not extracted for this review and are therefore not reported here. The authors do not state limitations in the abstract. The work shows the value of combining local and global representations and of post-processing, but it models context within the image rather than across slices, which is the dimension our project targets.

\subsection{Deep ensemble learning for joint segmentation and grading}
Wen \emph{et al.} (2025), in ``A deep ensemble learning framework for glioma segmentation and grading prediction'' \cite{wen2025}, observe that most methods treat glioma segmentation and risk-grade prediction as separate single tasks and that grading data are scarce. They propose a deep-ensemble framework built on U-Net that performs both tasks from pre-operative multi-modal MRI. The encoder uses asymmetric convolution and dual-domain attention to integrate information from the modalities, and a dual-branch decoder combines spatial features with global semantic information. A weighted composite adaptive loss balances the two objectives. The study is set on the BraTS dataset. The authors report that the method outperforms state-of-the-art approaches in both segmentation accuracy and grade prediction; exact values are not reported in the abstract. The architecture is convolutional with attention inside each image and does not model long-range dependencies between slices. It informs our project by showing that attention-augmented U-Nets remain a strong design for BraTS, while leaving inter-slice modelling open.

\subsection{SSL-MISS-Net: learning from incomplete sequences and annotations}
Song \emph{et al.} (2025), in ``An integrated MRI-based diagnostic framework for glioma with incomplete imaging sequences and imperfect annotations'' \cite{song2025}, address multi-centre glioma diagnosis, where differing imaging protocols and testing standards lead to missing MRI sequences and missing labels. They propose SSL-MISS-Net, which handles missing input sequences through cross-modal self-supervised learning and missing labels through a missing-label synergistic strategy. The study uses multi-centre cohorts from six in-house datasets and three public repositories (2,238 patients), evaluated with five-fold cross-validation and two independent test sets. Compared with the next-best method, the framework improved integrated-diagnosis accuracy by 4\% on validation and 10\% on test data and increased the amount of clinically usable data by 256\%. The task is diagnosis and molecular classification rather than voxel-level segmentation, so segmentation boundary quality is not applicable. The work is relevant because it shows the value of robustness to imperfect, real-world clinical data, a requirement our model would need to meet before clinical use.

\subsection{CrossFormer 2.5D model with radiomics for IDH/TERT prediction}
Niu \emph{et al.} (2025), in ``MRI transformer deep learning and radiomics for predicting IDH wild type TERT promoter mutant gliomas'' \cite{niu2025}, aim to identify IDH-wildtype, TERT-promoter-mutant gliomas from MRI and to stratify patient risk. They develop a multi-channel 2.5D deep-learning model and a 2D model, both based on the cross-scale attention vision transformer (CrossFormer), together with a radiomics model, and combine them in a stacking ensemble. Pre-operative T1CE and FLAIR images from 1,185 patients are analysed. The 2.5D model outperformed both the 2D and the radiomics models, with areas under the ROC curve (AUC) of 0.806--0.870, and the stacking model achieved the highest AUC of 0.855--0.904 across validation sets; Kaplan--Meier analysis showed significant survival differences between the predicted risk groups. The task is classification, so segmentation metrics are not applicable, and limitations are not reported in the abstract. Its finding that a 2.5D multi-slice input outperforms a 2D input on glioma MRI directly motivates our use of neighbouring-slice context for segmentation.

\subsection{Multi-task deep learning pipeline for glioma molecular status, grade and prognosis}
Wu \emph{et al.} (2024), in ``Biologically interpretable multi-task deep learning pipeline predicts molecular alterations, grade, and prognosis in glioma patients'' \cite{wu2024}, develop an end-to-end multi-task deep-learning pipeline that predicts IDH mutation, 1p/19q co-deletion and histological grade as auxiliary tasks and prognosis as the primary task, and they relate the predicted deep prognosis score to biological mechanisms. The study uses MRI from 2,776 patients: two private cohorts ($n = 1931$) for training and internal validation, and three public cohorts (TCGA, UCSF and EGD) for external validation. In external validation the pipeline reached AUCs of 0.892--0.903 for IDH status, 0.710--0.894 for 1p/19q status and 0.850--0.879 for grade, and C-indices of 0.723 (TCGA) and 0.671 (UCSF) for overall survival. This was the reference paper of our previous semester. Its pipeline depends on accurate tumour localisation, and the wide range of 1p/19q performance across cohorts illustrates how sensitive downstream analysis is to data quality. The work motivated our shift to the segmentation stage itself, since every downstream prediction depends on a reliable tumour mask.

\subsection{Cross-Slice Transformer with adaptive slice fusion (base paper)}
Abdikenov \emph{et al.} (2026), in ``Cross-slice transformer with adaptive slice fusion for lymph node segmentation in CT imaging'' \cite{abdikenov2026}, address lymph-node segmentation in CT, which is difficult because nodes are small, contrast is low and foreground voxels make up below 1\% of the volume. They propose the Cross-Slice Transformer (CST), a lightweight 2.5D architecture that takes five consecutive slices, extracts a statistical descriptor (mean, standard deviation, inter-slice difference and their interaction) from each slice, and applies transformer attention to these descriptors to learn adaptive slice-importance weights. The fused representation is decoded by a U-Net. The study uses two public TCIA datasets (296 and 206 patients) and compares CST with 2.5D U-Net, TransUNet and SwinUNet baselines and with a 3D nnU-Net. On the mediastinal dataset, CST achieved the highest Dice among 2.5D models ($0.543 \pm 0.044$) and the highest 2.5D recall (0.618) with only 7.97~M parameters, against 25--99~M for the 2.5D baselines, while nnU-Net-3D achieved the highest overall Dice (0.589). Ablations attribute gains of $+0.006$, $+0.008$ and $+0.008$ Dice to adaptive attention, the interaction descriptor and per-case intensity normalisation respectively. The limitations are its single-channel CT setting and its lower Dice than the 3D model. It is the base of our project: we adapt its central idea, attention across neighbouring slices, to four-channel MRI and to large nested glioma sub-regions, applying the attention to learned bottleneck features rather than to intensity statistics.

\subsection{U-Net}
Ronneberger \emph{et al.} (2015), in ``U-Net: convolutional networks for biomedical image segmentation'' \cite{ronneberger2015}, address the need for accurate segmentation from few annotated biomedical images. They propose the U-Net, a contracting path that captures context and a symmetric expanding path whose skip connections recover precise localisation, trained end to end with heavy elastic-deformation augmentation. The study uses electron-microscopy and light-microscopy cell-segmentation challenges from the ISBI workshops. The network won the ISBI 2015 cell-tracking challenge, with an intersection-over-union of 0.9203 on the PhC-U373 dataset and 0.7756 on DIC-HeLa, and it segments a $512 \times 512$ image in under a second on a GPU. The original work is 2D and does not address volumetric consistency. The U-Net is the backbone of both our proposed model and our baseline.

\subsection{3D U-Net}
\c{C}i\c{c}ek \emph{et al.} (2016), in ``3D U-Net: learning dense volumetric segmentation from sparse annotation'' \cite{cicek2016}, extend the U-Net to volumetric data by replacing 2D operations with 3D convolutions, pooling and up-convolutions, and they show that dense 3D segmentation can be learned from a few annotated 2D slices. The study uses confocal microscopy volumes of the \emph{Xenopus} kidney and covers both a semi-automatic setting (densifying sparse annotations) and a fully automatic setting (segmenting unseen volumes); the authors report intersection-over-union and show that the 3D network outperforms slice-wise 2D processing. The main limitation for clinical MRI is the memory cost of 3D feature maps, which constrains input size and batch size. The work establishes the benefit of inter-slice context, which our 2.5D model seeks at lower cost.

\subsection{nnU-Net}
Isensee \emph{et al.} (2021), in ``nnU-Net: a self-configuring method for deep learning-based biomedical image segmentation'' \cite{isensee2021}, argue that designing segmentation solutions is highly dependent on dataset properties and hardware. They propose nnU-Net, which automatically configures preprocessing, network architecture, training and post-processing for a new task from a set of fixed parameters, interdependent rules and empirical decisions. Without manual intervention, nnU-Net surpassed most existing approaches, including highly specialised solutions, on 23 public datasets used in international segmentation competitions. Its limitation in our setting is computational: its full-resolution 3D configuration and long training schedule exceed the resources available to us for most of the project. It informs our work as the strongest standard reference and as evidence that careful training protocol and evaluation matter as much as architecture.

\subsection{TransBTS}
Wang \emph{et al.} (2021), in ``TransBTS: multimodal brain tumor segmentation using transformer'' \cite{wang2021transbts}, aim to add long-range dependency modelling to 3D CNNs for brain-tumour segmentation. They encode the multi-modal MRI volume with a 3D CNN, model global context at the bottleneck with a transformer applied to the resulting tokens, and decode with a 3D CNN. The study uses the BraTS 2019 and BraTS 2020 datasets, evaluated with Dice and Hausdorff distance on the official validation sets. On the BraTS 2020 validation set the authors report Dice scores of 0.7873 (ET), 0.9009 (WT) and 0.8173 (TC). The method relies on 3D patches and their memory cost. Its idea of placing a transformer at the U-Net bottleneck is shared by our design, but our transformer attends across slices rather than across spatial positions of a 3D volume.

\subsection{Swin UNETR}
Hatamizadeh \emph{et al.} (2022), in ``Swin UNETR: Swin transformers for semantic segmentation of brain tumors in MRI images'' \cite{hatamizadeh2022swin}, extend their earlier UNETR \cite{hatamizadeh2022unetr} by using a Swin transformer encoder with shifted-window self-attention to extract features at five resolutions from multi-modal 3D MRI, connected to a convolutional decoder through skip connections. The study is set in the BraTS 2021 challenge, evaluated with Dice and HD95, and the method was reported among the top-ranking approaches in the challenge validation phase. Like other 3D transformers it requires substantial GPU memory and training time. It represents the state of transformer-based 3D brain-tumour segmentation and indicates the accuracy that full volumetric context can reach when resources allow.

\clearpage\begin{landscape}
\section{Comparison of reviewed literature}
Table~\ref{tab:literature} compares the twelve reviewed works. Study context covers the data and setting, and missing details are marked ``Not reported''.

\begingroup\small\setlength{\tabcolsep}{4pt}\renewcommand{\arraystretch}{1.2}
\begin{longtable}{|p{2.2cm}|p{2.6cm}|p{3.6cm}|p{2.7cm}|p{2.6cm}|p{3.2cm}|p{3.2cm}|}
\caption{Comparison of reviewed literature}\label{tab:literature}\\\hline
\textbf{Work and citation} & \textbf{Problem or objective} & \textbf{Approach and contribution} & \textbf{Study context} & \textbf{Evaluation criteria} & \textbf{Key findings} & \textbf{Limitations and gap}\\\hline\endfirsthead
\textbf{Work and citation} & \textbf{Problem or objective} & \textbf{Approach and contribution} & \textbf{Study context} & \textbf{Evaluation criteria} & \textbf{Key findings} & \textbf{Limitations and gap}\\\hline\endhead
Work 01 \cite{aslam2025} Aslam, 2025 & Brain-tumour segmentation from multi-modal MRI & Multi-scale contextual aggregation and gated attention fusion (MM-MSCA-AF) & BraTS 2020; T1, T2, FLAIR, T1ce & Dice & Dice 0.8589 overall, 0.8158 necrotic region; better than U-Net, nnU-Net, Attention U-Net & No inter-slice modelling; HD95 not reported\\\hline
Work 02 \cite{li2025deepglioseg} Li, 2025 & Glioma segmentation under heterogeneity and region imbalance & CNN--transformer parallel module, region-probability weighting, TTA and volume-constrained post-processing & Three public datasets and one private clinical dataset & Quantitative and qualitative segmentation metrics & Reported superior to state-of-the-art methods; values not reported here & Context within images, not across slices\\\hline
Work 03 \cite{wen2025} Wen, 2025 & Joint glioma segmentation and risk grading & U-Net ensemble with asymmetric convolution, dual-domain attention, dual-branch decoder, adaptive composite loss & BraTS & Segmentation accuracy; grading accuracy & Reported better than state-of-the-art in both tasks; values not reported in abstract & Convolutional; no long-range inter-slice dependency\\\hline
Work 04 \cite{song2025} Song, 2025 & Diagnosis with missing sequences and labels & SSL-MISS-Net: cross-modal self-supervision and missing-label strategy & 2,238 patients, nine multi-centre cohorts; five-fold CV and two test sets & Diagnostic accuracy & +4\% (validation) and +10\% (test) accuracy over next-best; +256\% usable data & Classification, not voxel-level segmentation\\\hline
Work 05 \cite{niu2025} Niu, 2025 & Predict IDH-wt, TERTp-mutant gliomas; risk stratification & CrossFormer 2.5D and 2D models plus radiomics in a stacking ensemble & 1,185 patients; T1CE and FLAIR & AUC; Kaplan--Meier, log-rank & 2.5D AUC 0.806--0.870 beat 2D and radiomics; stacking AUC 0.855--0.904 & Classification only; supports value of 2.5D input\\\hline
Work 06 \cite{wu2024} Wu, 2024 & Predict molecular status, grade and prognosis & End-to-end multi-task deep learning with interpretable prognosis score & 2,776 patients; private training, three public external cohorts & AUC; C-index & External AUC 0.892--0.903 (IDH), 0.710--0.894 (1p/19q), 0.850--0.879 (grade); C-index 0.723 / 0.671 & Depends on tumour localisation; segmentation not the focus\\\hline
Work 07 \cite{abdikenov2026} Abdikenov, 2026 & Lymph-node segmentation in CT & 2.5D Cross-Slice Transformer: per-slice descriptors with transformer attention for adaptive slice fusion & Two TCIA CT datasets (296 and 206 patients) & Dice, recall, node detection rate, parameters & Best 2.5D Dice $0.543 \pm 0.044$ with 7.97~M parameters; nnU-Net-3D 0.589 & Single-channel CT, small targets; untested on multi-modal MRI\\\hline
Work 08 \cite{ronneberger2015} Ronneberger, 2015 & Segmentation from few annotations & U-Net encoder--decoder with skip connections and elastic augmentation & ISBI EM and cell-tracking challenges & IoU; warping error & Won ISBI 2015 cell tracking: IoU 0.9203 (PhC-U373), 0.7756 (DIC-HeLa) & 2D only; no volumetric consistency\\\hline
Work 09 \cite{cicek2016} \c{C}i\c{c}ek, 2016 & Volumetric segmentation from sparse labels & 3D U-Net with 3D convolutions; learns from sparsely annotated slices & \emph{Xenopus} kidney confocal volumes & IoU & 3D network outperformed slice-wise 2D processing & High memory cost of 3D feature maps\\\hline
Work 10 \cite{isensee2021} Isensee, 2021 & Dataset-specific design effort & Self-configuring nnU-Net pipeline & 23 public segmentation datasets & Challenge metrics (Dice-based) & Surpassed most existing, including specialised, methods & Heavy 3D training; resource-intensive\\\hline
Work 11 \cite{wang2021transbts} Wang, 2021 & Long-range context for 3D brain-tumour segmentation & 3D CNN encoder, transformer bottleneck, 3D CNN decoder & BraTS 2019 and 2020 official validation & Dice; Hausdorff distance & BraTS 2020 validation Dice: ET 0.7873, WT 0.9009, TC 0.8173 & 3D memory cost; spatial (not slice-wise) attention\\\hline
Work 12 \cite{hatamizadeh2022swin} Hatamizadeh, 2022 & Brain-tumour segmentation with hierarchical transformers & Swin transformer encoder at five resolutions with CNN decoder & BraTS 2021 challenge & Dice; HD95 & Among top-ranking methods in validation phase & High GPU memory and training time\\\hline
\end{longtable}\endgroup\end{landscape}
